\addcontentsline{toc}{chapter}{Resumen}

\chapter*{Resumen}

\vspace*{-8mm}

Se presenta una implementación de Redes Neuronales Informadas por Física con
activación sinusoidal (PINN-SIREN) para la simulación acelerada del patrón de
speckle óptico, mediante la resolución de la ecuación de Helmholtz bidimensional
con campo eléctrico complejo y condición de frontera de fase aleatoria.
La arquitectura SIREN propuesta ($5 \times 128$ neuronas, $\omega_0 = 1.0$)
alcanza errores $L^2$ de $0.006\%$ en el caso 1D y $0.171\%$ en el caso 2D
respecto a soluciones analíticas, varios órdenes de magnitud por debajo del
umbral de tesis ($5\%$).
El patrón de speckle se genera imponiendo una fase aleatoria
$\phi(x) \sim \mathcal{U}(0, 2\pi)$ en la superficie iluminada y resolviendo la
ecuación de Helmholtz libremente en el interior del dominio.
La validación estadística sigue el criterio de Goodman: contraste
$C = \sigma_I / \langle I \rangle = 1.0253$ y distribución de intensidades
exponencial negativa, confirmando speckle completamente desarrollado.
Una estrategia de optimización bifásica Adam~+~L-BFGS y el muestreo por
hipercubo latino (LHS) son clave para la convergencia.
La reproducibilidad del método fue verificada con tres semillas independientes,
obteniendo un error $L^2$ promedio de $0.155 \pm 0.020\%$.

\bigskip

\noindent\textbf{Palabras clave:} Redes Neuronales Informadas por Física,
SIREN, ecuación de Helmholtz, speckle óptico, activación sinusoidal,
muestreo por hipercubo latino.
